% Options for packages loaded elsewhere
% Options for packages loaded elsewhere
\PassOptionsToPackage{unicode}{hyperref}
\PassOptionsToPackage{hyphens}{url}
\PassOptionsToPackage{dvipsnames,svgnames,x11names}{xcolor}
%
\documentclass[
  letterpaper,
  DIV=11,
  numbers=noendperiod]{scrartcl}
\usepackage{xcolor}
\usepackage{amsmath,amssymb}
\setcounter{secnumdepth}{5}
\usepackage{iftex}
\ifPDFTeX
  \usepackage[T1]{fontenc}
  \usepackage[utf8]{inputenc}
  \usepackage{textcomp} % provide euro and other symbols
\else % if luatex or xetex
  \usepackage{unicode-math} % this also loads fontspec
  \defaultfontfeatures{Scale=MatchLowercase}
  \defaultfontfeatures[\rmfamily]{Ligatures=TeX,Scale=1}
\fi
\usepackage{lmodern}
\ifPDFTeX\else
  % xetex/luatex font selection
\fi
% Use upquote if available, for straight quotes in verbatim environments
\IfFileExists{upquote.sty}{\usepackage{upquote}}{}
\IfFileExists{microtype.sty}{% use microtype if available
  \usepackage[]{microtype}
  \UseMicrotypeSet[protrusion]{basicmath} % disable protrusion for tt fonts
}{}
\makeatletter
\@ifundefined{KOMAClassName}{% if non-KOMA class
  \IfFileExists{parskip.sty}{%
    \usepackage{parskip}
  }{% else
    \setlength{\parindent}{0pt}
    \setlength{\parskip}{6pt plus 2pt minus 1pt}}
}{% if KOMA class
  \KOMAoptions{parskip=half}}
\makeatother
% Make \paragraph and \subparagraph free-standing
\makeatletter
\ifx\paragraph\undefined\else
  \let\oldparagraph\paragraph
  \renewcommand{\paragraph}{
    \@ifstar
      \xxxParagraphStar
      \xxxParagraphNoStar
  }
  \newcommand{\xxxParagraphStar}[1]{\oldparagraph*{#1}\mbox{}}
  \newcommand{\xxxParagraphNoStar}[1]{\oldparagraph{#1}\mbox{}}
\fi
\ifx\subparagraph\undefined\else
  \let\oldsubparagraph\subparagraph
  \renewcommand{\subparagraph}{
    \@ifstar
      \xxxSubParagraphStar
      \xxxSubParagraphNoStar
  }
  \newcommand{\xxxSubParagraphStar}[1]{\oldsubparagraph*{#1}\mbox{}}
  \newcommand{\xxxSubParagraphNoStar}[1]{\oldsubparagraph{#1}\mbox{}}
\fi
\makeatother

\usepackage{color}
\usepackage{fancyvrb}
\newcommand{\VerbBar}{|}
\newcommand{\VERB}{\Verb[commandchars=\\\{\}]}
\DefineVerbatimEnvironment{Highlighting}{Verbatim}{commandchars=\\\{\}}
% Add ',fontsize=\small' for more characters per line
\usepackage{framed}
\definecolor{shadecolor}{RGB}{241,243,245}
\newenvironment{Shaded}{\begin{snugshade}}{\end{snugshade}}
\newcommand{\AlertTok}[1]{\textcolor[rgb]{0.68,0.00,0.00}{#1}}
\newcommand{\AnnotationTok}[1]{\textcolor[rgb]{0.37,0.37,0.37}{#1}}
\newcommand{\AttributeTok}[1]{\textcolor[rgb]{0.40,0.45,0.13}{#1}}
\newcommand{\BaseNTok}[1]{\textcolor[rgb]{0.68,0.00,0.00}{#1}}
\newcommand{\BuiltInTok}[1]{\textcolor[rgb]{0.00,0.23,0.31}{#1}}
\newcommand{\CharTok}[1]{\textcolor[rgb]{0.13,0.47,0.30}{#1}}
\newcommand{\CommentTok}[1]{\textcolor[rgb]{0.37,0.37,0.37}{#1}}
\newcommand{\CommentVarTok}[1]{\textcolor[rgb]{0.37,0.37,0.37}{\textit{#1}}}
\newcommand{\ConstantTok}[1]{\textcolor[rgb]{0.56,0.35,0.01}{#1}}
\newcommand{\ControlFlowTok}[1]{\textcolor[rgb]{0.00,0.23,0.31}{\textbf{#1}}}
\newcommand{\DataTypeTok}[1]{\textcolor[rgb]{0.68,0.00,0.00}{#1}}
\newcommand{\DecValTok}[1]{\textcolor[rgb]{0.68,0.00,0.00}{#1}}
\newcommand{\DocumentationTok}[1]{\textcolor[rgb]{0.37,0.37,0.37}{\textit{#1}}}
\newcommand{\ErrorTok}[1]{\textcolor[rgb]{0.68,0.00,0.00}{#1}}
\newcommand{\ExtensionTok}[1]{\textcolor[rgb]{0.00,0.23,0.31}{#1}}
\newcommand{\FloatTok}[1]{\textcolor[rgb]{0.68,0.00,0.00}{#1}}
\newcommand{\FunctionTok}[1]{\textcolor[rgb]{0.28,0.35,0.67}{#1}}
\newcommand{\ImportTok}[1]{\textcolor[rgb]{0.00,0.46,0.62}{#1}}
\newcommand{\InformationTok}[1]{\textcolor[rgb]{0.37,0.37,0.37}{#1}}
\newcommand{\KeywordTok}[1]{\textcolor[rgb]{0.00,0.23,0.31}{\textbf{#1}}}
\newcommand{\NormalTok}[1]{\textcolor[rgb]{0.00,0.23,0.31}{#1}}
\newcommand{\OperatorTok}[1]{\textcolor[rgb]{0.37,0.37,0.37}{#1}}
\newcommand{\OtherTok}[1]{\textcolor[rgb]{0.00,0.23,0.31}{#1}}
\newcommand{\PreprocessorTok}[1]{\textcolor[rgb]{0.68,0.00,0.00}{#1}}
\newcommand{\RegionMarkerTok}[1]{\textcolor[rgb]{0.00,0.23,0.31}{#1}}
\newcommand{\SpecialCharTok}[1]{\textcolor[rgb]{0.37,0.37,0.37}{#1}}
\newcommand{\SpecialStringTok}[1]{\textcolor[rgb]{0.13,0.47,0.30}{#1}}
\newcommand{\StringTok}[1]{\textcolor[rgb]{0.13,0.47,0.30}{#1}}
\newcommand{\VariableTok}[1]{\textcolor[rgb]{0.07,0.07,0.07}{#1}}
\newcommand{\VerbatimStringTok}[1]{\textcolor[rgb]{0.13,0.47,0.30}{#1}}
\newcommand{\WarningTok}[1]{\textcolor[rgb]{0.37,0.37,0.37}{\textit{#1}}}

\usepackage{longtable,booktabs,array}
\newcounter{none} % for unnumbered tables
\usepackage{calc} % for calculating minipage widths
% Correct order of tables after \paragraph or \subparagraph
\usepackage{etoolbox}
\makeatletter
\patchcmd\longtable{\par}{\if@noskipsec\mbox{}\fi\par}{}{}
\makeatother
% Allow footnotes in longtable head/foot
\IfFileExists{footnotehyper.sty}{\usepackage{footnotehyper}}{\usepackage{footnote}}
\makesavenoteenv{longtable}
\usepackage{graphicx}
\makeatletter
\newsavebox\pandoc@box
\newcommand*\pandocbounded[1]{% scales image to fit in text height/width
  \sbox\pandoc@box{#1}%
  \Gscale@div\@tempa{\textheight}{\dimexpr\ht\pandoc@box+\dp\pandoc@box\relax}%
  \Gscale@div\@tempb{\linewidth}{\wd\pandoc@box}%
  \ifdim\@tempb\p@<\@tempa\p@\let\@tempa\@tempb\fi% select the smaller of both
  \ifdim\@tempa\p@<\p@\scalebox{\@tempa}{\usebox\pandoc@box}%
  \else\usebox{\pandoc@box}%
  \fi%
}
% Set default figure placement to htbp
\def\fps@figure{htbp}
\makeatother





\setlength{\emergencystretch}{3em} % prevent overfull lines

\providecommand{\tightlist}{%
  \setlength{\itemsep}{0pt}\setlength{\parskip}{0pt}}



 


\KOMAoption{captions}{tableheading}
\makeatletter
\@ifpackageloaded{caption}{}{\usepackage{caption}}
\AtBeginDocument{%
\ifdefined\contentsname
  \renewcommand*\contentsname{Table of contents}
\else
  \newcommand\contentsname{Table of contents}
\fi
\ifdefined\listfigurename
  \renewcommand*\listfigurename{List of Figures}
\else
  \newcommand\listfigurename{List of Figures}
\fi
\ifdefined\listtablename
  \renewcommand*\listtablename{List of Tables}
\else
  \newcommand\listtablename{List of Tables}
\fi
\ifdefined\figurename
  \renewcommand*\figurename{Figure}
\else
  \newcommand\figurename{Figure}
\fi
\ifdefined\tablename
  \renewcommand*\tablename{Table}
\else
  \newcommand\tablename{Table}
\fi
}
\@ifpackageloaded{float}{}{\usepackage{float}}
\floatstyle{ruled}
\@ifundefined{c@chapter}{\newfloat{codelisting}{h}{lop}}{\newfloat{codelisting}{h}{lop}[chapter]}
\floatname{codelisting}{Listing}
\newcommand*\listoflistings{\listof{codelisting}{List of Listings}}
\makeatother
\makeatletter
\makeatother
\makeatletter
\@ifpackageloaded{caption}{}{\usepackage{caption}}
\@ifpackageloaded{subcaption}{}{\usepackage{subcaption}}
\makeatother
\usepackage{bookmark}
\IfFileExists{xurl.sty}{\usepackage{xurl}}{} % add URL line breaks if available
\urlstyle{same}
% fallback for those not using the hyperref driver hyperxmp:
\makeatletter
\@ifundefined{xmpquote}{\newcommand{\xmpquote}[1]{#1}}{}
\makeatother
\hypersetup{
  pdftitle={ECOM6004 Assessment 1 - Binary and Ordered Response Models},
  pdfauthor={Student ID: 12345678},
  colorlinks=true,
  linkcolor={blue},
  filecolor={Maroon},
  citecolor={Blue},
  urlcolor={Blue},
  pdfcreator={LaTeX via pandoc}}


\title{ECOM6004 Assessment 1 - Binary and Ordered Response Models}
\author{Student ID: 12345678}
\date{Invalid Date}
\begin{document}
\maketitle

\renewcommand*\contentsname{Table of contents}
{
\setcounter{tocdepth}{2}
\tableofcontents
}

\section{Introduction}\label{introduction}

This report presents the analysis of the Bank Marketing dataset using
binary and ordered response models. The analysis follows the ECOM6004
Assessment 1 requirements for Semester 2, 2026.

\section{Data Description}\label{data-description}

The Bank Marketing dataset contains \texttt{\{\ len(full\_df)\ \}}
observations with the following key variables:

\begin{itemize}
\tightlist
\item
  \textbf{age}: Age of the client
\item
  \textbf{balance}: Average yearly balance in euros
\item
  \textbf{housing}: Whether the client has a housing loan
\item
  \textbf{loan}: Whether the client has a personal loan
\item
  \textbf{contact}: Contact communication type
\item
  \textbf{previous}: Number of contacts performed before this campaign
\item
  \textbf{y}: Target variable - whether the client subscribed to a term
  deposit
\end{itemize}

The dataset was split into training (80\%) and test (20\%) sets using
stratified sampling with the student ID as the random seed.

\section{Binary Response Models}\label{binary-response-models}

\subsection{Model Specification}\label{model-specification}

Three binary response models were estimated on the training data:

\begin{enumerate}
\def\labelenumi{\arabic{enumi}.}
\tightlist
\item
  Linear Probability Model (LPM)
\item
  Logit Model
\item
  Probit Model
\end{enumerate}

All models use the following specification:

\texttt{target\ \textasciitilde{}\ age10\ +\ balance1000\ +\ C(housing)\ +\ C(loan)\ +\ C(contact)\ +\ previous\ +\ C(prior\_status)}

Where: - \texttt{age10} = (age - 40) / 10 - \texttt{balance1000} =
balance / 1000 - \texttt{prior\_status} is derived from \texttt{pdays}
and \texttt{poutcome}

\subsection{Model Results}\label{model-results}

\begin{Shaded}
\begin{Highlighting}[]
\InformationTok{\textasciigrave{}\textasciigrave{}\textasciigrave{}\{python\}}
\CommentTok{\#| echo: false}
\CommentTok{\#| warning: false}

\ImportTok{from}\NormalTok{ src.data\_loader }\ImportTok{import}\NormalTok{ load\_and\_prepare\_bank\_data}
\ImportTok{from}\NormalTok{ src.binary\_models }\ImportTok{import}\NormalTok{ fit\_binary\_models}
\ImportTok{from}\NormalTok{ src.utils }\ImportTok{import}\NormalTok{ format\_summary\_table}

\NormalTok{STUDENT\_ID }\OperatorTok{=} \DecValTok{12345678}
\NormalTok{full\_df, train\_df, test\_df }\OperatorTok{=}\NormalTok{ load\_and\_prepare\_bank\_data(}\StringTok{"data/bank\_assessment\_sem2\_2026.csv"}\NormalTok{, STUDENT\_ID)}
\NormalTok{binary\_results }\OperatorTok{=}\NormalTok{ fit\_binary\_models(train\_df)}
\NormalTok{coef\_table }\OperatorTok{=}\NormalTok{ format\_summary\_table(binary\_results, [}\StringTok{\textquotesingle{}lpm\textquotesingle{}}\NormalTok{, }\StringTok{\textquotesingle{}logit\textquotesingle{}}\NormalTok{, }\StringTok{\textquotesingle{}probit\textquotesingle{}}\NormalTok{])}
\NormalTok{coef\_table}
\InformationTok{\textasciigrave{}\textasciigrave{}\textasciigrave{}}
\end{Highlighting}
\end{Shaded}

\subsection{Model Comparison}\label{model-comparison}

The table above compares the coefficients across the three models. Key
observations include:

\begin{itemize}
\tightlist
\item
  {[}Interpretation of age10 coefficient{]}
\item
  {[}Interpretation of balance1000 coefficient{]}
\item
  {[}Interpretation of housing coefficient{]}
\item
  {[}Interpretation of loan coefficient{]}
\item
  {[}Interpretation of contact coefficient{]}
\item
  {[}Interpretation of previous coefficient{]}
\item
  {[}Interpretation of prior\_status coefficients{]}
\end{itemize}

\subsection{Out-of-Sample Evaluation}\label{out-of-sample-evaluation}

The logit model was evaluated on the held-out test set using a 0.50
threshold.

\begin{Shaded}
\begin{Highlighting}[]
\InformationTok{\textasciigrave{}\textasciigrave{}\textasciigrave{}\{python\}}
\CommentTok{\#| echo: false}
\CommentTok{\#| warning: false}

\ImportTok{from}\NormalTok{ src.binary\_models }\ImportTok{import}\NormalTok{ evaluate\_test\_set}
\ImportTok{from}\NormalTok{ src.utils }\ImportTok{import}\NormalTok{ format\_metrics\_table}

\NormalTok{test\_metrics }\OperatorTok{=}\NormalTok{ evaluate\_test\_set(binary\_results[}\StringTok{\textquotesingle{}logit\textquotesingle{}}\NormalTok{], test\_df)}
\NormalTok{metrics\_table }\OperatorTok{=}\NormalTok{ format\_metrics\_table(test\_metrics)}
\NormalTok{metrics\_table}
\InformationTok{\textasciigrave{}\textasciigrave{}\textasciigrave{}}
\end{Highlighting}
\end{Shaded}

The test set evaluation shows: - Observed event rate:
\texttt{\{\ test\_metrics{[}\textquotesingle{}obs\_event\_rate\textquotesingle{}{]}:.4f\ \}}
- Mean predicted probability:
\texttt{\{\ test\_metrics{[}\textquotesingle{}mean\_pred\_prob\textquotesingle{}{]}:.4f\ \}}
- Brier score:
\texttt{\{\ test\_metrics{[}\textquotesingle{}brier\_score\textquotesingle{}{]}:.4f\ \}}
- Sensitivity:
\texttt{\{\ test\_metrics{[}\textquotesingle{}sensitivity\textquotesingle{}{]}:.4f\ \}}
- Specificity:
\texttt{\{\ test\_metrics{[}\textquotesingle{}specificity\textquotesingle{}{]}:.4f\ \}}
- Balanced accuracy:
\texttt{\{\ test\_metrics{[}\textquotesingle{}balanced\_accuracy\textquotesingle{}{]}:.4f\ \}}

\section{Ordered Response Models}\label{ordered-response-models}

\subsection{Model Specification}\label{model-specification-1}

Ordered Logit and Ordered Probit models were estimated on the full
dataset using the \texttt{response\_level} variable as the ordered
outcome.

\subsection{Model Results}\label{model-results-1}

\begin{Shaded}
\begin{Highlighting}[]
\InformationTok{\textasciigrave{}\textasciigrave{}\textasciigrave{}\{python\}}
\CommentTok{\#| echo: false}
\CommentTok{\#| warning: false}

\ImportTok{from}\NormalTok{ src.ordered\_models }\ImportTok{import}\NormalTok{ fit\_ordered\_models}

\NormalTok{ordered\_results }\OperatorTok{=}\NormalTok{ fit\_ordered\_models(full\_df)}
\NormalTok{ordered\_coef\_table }\OperatorTok{=}\NormalTok{ format\_summary\_table(ordered\_results, [}\StringTok{\textquotesingle{}ologit\textquotesingle{}}\NormalTok{, }\StringTok{\textquotesingle{}oprobit\textquotesingle{}}\NormalTok{])}
\NormalTok{ordered\_coef\_table}
\InformationTok{\textasciigrave{}\textasciigrave{}\textasciigrave{}}
\end{Highlighting}
\end{Shaded}

\subsection{Interpretation}\label{interpretation}

The ordered model coefficients represent the effect of each predictor on
the log-odds of being in a higher response category. Key findings
include:

\begin{itemize}
\tightlist
\item
  {[}Interpretation of key coefficients{]}
\item
  {[}Comparison between ordered logit and ordered probit{]}
\end{itemize}

\section{Conclusion}\label{conclusion}

This analysis demonstrates the application of binary and ordered
response models to the Bank Marketing dataset. The binary models show
{[}summary of findings{]}, while the ordered models provide insights
into {[}summary of findings{]}.

\section{References}\label{references}

{[}Add any references here if applicable{]}




\end{document}
